\chapter*{Appendices}
\addcontentsline{toc}{chapter}{Appendices}

\section*{Appendix 1 -- To Be an Authority}

Excerpts from \emph{The Retreat of the State: The Diffusion of Power in the World Economy} by Susan Strange, 1996.

\small
\begin{longtable}{p{1.2cm} p{13cm}}
\toprule
\textbf{Page} & \textbf{Quote} \\
\midrule
\endhead

p.xiii & The implicit assumption conveyed by the two words, `global' and `governance', is that government is being achieved on a world scale by a world \textbf{authority}. \\
\midrule

p.xv & I firmly believe that the new realism of the Stopford-Strange analysis of corporate strategies and state development policies makes it imperative to look seriously at the power exercised by \textbf{authorities} other than states. \\
\midrule

p.115 & Representatives of the Italian state in Rome began to think better of their delegation of \textbf{authority} to the mafia as the latter tended to act more and more blatantly outside the law. A new generation of well-trained magistrates, themselves mostly Sicilians, were courageous enough to initiate an open break between state and mafia. \\
\midrule

p.116 & Sociologists have argued that criminal gangs, like underground resistance movements in wartime or recalcitrant groups in prisons, tend to emerge when state \textbf{authority}, for whatever reason is already weakened, and the government has lost or failed to obtain the consent of the governed. \\
\midrule

p.165 & Here, we must distinguish between functions and \textbf{authority} that are consciously delegated to the organisation by member states, and functions and \textbf{authority} that have been assumed by the officials, independently of the wishes or decisions of member states. \\
\midrule

p.168 & The net result of leaving each creditor to define the terms and modalities of debt relief was that \textbf{authority} was delegated on a case-by-case basis to international institutions other than the United Nations or the IMF or IBRD --- the so-called Paris Club or official creditors and the London Club of private creditors. There are other less politically important matters on which member states have delegated executive \textbf{authority} --- or seemed to do so --- to international institutions. \\
\midrule

p.171 & Four major examples show that such delegation is conceded only when the system --- some part of the structures of the world market economy --- is perceived as being seriously at risk, and when the direct and indirect costs of delegated \textbf{authority} are relatively insignificant. Here we may ignore the Council of Ministers, as being clearly an inter-governmental body whose members are named by and responsible to national governments. It is only important insofar as it delegates \textbf{authority} to the Commission, to the Parliament or to the European Court, independent of national governments. \\
\midrule

p.173 & Whether this trend in Europe is peculiar to the European Community or is a harbinger of a more general shift of \textbf{authority} to international judicial bodies is a much more open question. \\
\midrule

p.174 & Although German resistance on behalf of national regulatory \textbf{authority} stopped the complete transfer of \textbf{authority} over mergers to the EC, some significant shift did take place. \\
\midrule

p.179 & Transfer of \textbf{authority} from the member governments to federal institutions over this central responsibility of political \textbf{authority} in a market economy has not yet happened. \\
\midrule

p.187 & Another feature which the triangular model also accommodates is the fact that there are striking variations across sectors in the nature and kind of \textbf{authority} and how much it, or they, intervene with the play of market forces. \\
\midrule

p.189 & Mine's argument in \emph{Le Nouveau Moyen Age} is that these areas without legitimate, acknowledged \textbf{authority}, in which the law of the jungle rules, are growing, especially in Africa and in the former Soviet Union. \textbf{authority} is divided between the formal institutions of the state and local potentates, chiefs or gang leaders; between vassal and suzerain, the responsibility for keeping order is as unclear as it was in the middle ages. \\
\midrule

p.192 & Although there were occasions when the delegation of \textbf{authority} to an international institution, as to any other body, seemed to give it some independent power of its own, that was usually more an illusion than reality. \\
\midrule

p.198 & To make \textbf{authority} acceptable, effective and respected, there has to be some combination of forces to check the arbitrary or self-serving use of power and to see that it is used at least in part for the common good. \\
\bottomrule
\end{longtable}

\noindent Source: STRANGE, 1996.

\normalsize

\section*{Appendix 2 -- To Exercise Authority}

Excerpts from \emph{The Retreat of the State: The Diffusion of Power in the World Economy} by Susan Strange, 1996.

\small
\begin{longtable}{p{1.2cm} p{13cm}}
\toprule
\textbf{Page} & \textbf{Quote} \\
\midrule
\endhead

p.110 & For example, its \textbf{authority} --- like that of a state --- is exercised through an established power structure, by means of which obedience is rewarded and disobedience punished, occasionally by the use of violence and always by the threat of violence. \\
\midrule

p.133 & {[}\ldots{]} \textbf{authority} in political economy is recognisable by the power to alter or modify the behaviour of others by using incentives and disincentives to affect the choice and range of options, {[}\ldots{]} \\
\midrule

p.171 & The conclusion must be that IOs, both in their dependent and independent exercise of \textbf{authority}, are essentially system-preserving. \\
\midrule

p.184 & The first, basic question was `Who, or what, is responsible for change?' The second was `Who, or what, exercises \textbf{authority} --- the power to alter \textbf{outcomes} and redefine options for others --- in the world economy or world society?' \\
\midrule

p.196 & The only other important consequences of the retreat of the state and the diffusion of state \textbf{authority} sketched in earlier chapters relates to legitimacy and democracy. \\
\midrule

p.197 & But if those institutions are now suffering the kind of diffusion of \textbf{authority} I have described, not much remains of the accountability of market forces to political constraints. Moreover, none of the non-state \textbf{authorities} to whom \textbf{authority} has shifted, is democratically governed. Firms --- the new players in transnational economic diplomacy --- are hierarchies, not democracies. \\
\midrule

p.199 & With the end of the Cold War, and with the triumph of the market economy, there is a new absence of absolutes. In a world of multiple, diffused \textbf{authority}, each of us shares Pinocchio's problem; our individual consciences are our only guide. \\
\bottomrule
\end{longtable}

\noindent Source: STRANGE, 1996.

\normalsize

\section*{Appendix 3 -- Control, 1988}

Excerpts on \textbf{control} in \emph{States and Markets} by Susan Strange, 1988.

\small
\begin{longtable}{p{1.2cm} p{9.5cm} p{3.2cm}}
\toprule
\textbf{Page} & \textbf{Quote} & \textbf{Object} \\
\midrule
\endhead

p.30 & But finance --- the \textbf{control} of credit --- is the facet which has perhaps risen in importance in the last quarter century more rapidly than any other and has come to be of decisive importance in international economic relations and in the competition of corporate enterprises. & Credit \\
\midrule

p.32--33 & But the power of the ayatollahs in defending and promoting Islamic virtues would have been constrained if they had not also gained \textbf{control} over the state and the armed forces sufficient to confirm their \textbf{authority} both within the country and beyond. & State and armed forces \\
\midrule

p.37 & Too often, they have ignored or refused to contemplate structural power, or the power to define the structure, to choose the game as well as to set the rules under which it is to be played. It is as if you said, `This man has power in relation to this woman because he can knock her down', ignoring the fact of structural power in a masculine-dominated social structure that gives the man social status, legal rights and \textbf{control} over the family money that makes it unnecessary even to threaten to knock her down unless she does as she is told. & Family money \\
\midrule

p.84 & The justice or injustice of the distributional effects of change in the production structure in short has been uneven, complex and subjective. Some wider issues concerning the system in general remain to be considered. Among these, there are three, of which the first --- whether states have the power to \textbf{control} the transnational corporations --- is familiar to most people and has been much discussed. & Transnational corporations \\
\midrule

p.85 & At the national level, the developing states that wanted to have the mandatory `shall' instead of the advisory `should' put into a UN Code were meanwhile seeking to use national political power against the foreign corporations. {[}\ldots{]} the displaced companies kept \textbf{control} over market access, by making long-term contracts with the customers, for instance. They also had command of the technology necessary to remain competitive in world markets. & Market access \\
\midrule

p.86 & The question now is whether this traditionally exclusive power claimed by all states alike is being eroded by large corporations through their \textbf{control} over the production structure. & Production structure \\
\midrule

p.101 & And there was one final factor that allowed the steady outflow of British capital before 1914 to help the world economy to grow reasonably steadily. It was India. {[}\ldots{]} But Britain's political \textbf{control} over India allowed London to extract annual shipments of gold in respect of Home Charges and to use its \textbf{control} over the sterling-rupee exchange rate and other devices to stop the gold seeping back to the Indian economy. & India / Sterling-rupee exchange rate \\
\midrule

p.106 & Briefly, the Eurodollar market (and later markets for Eurosterling, Euromarks, Euroyen, etc.) developed because of two inviting gaps in government \textbf{controls} over the power of banks to create credit. & Power of banks to create credit \\
\midrule

p.123 & In the knowledge structure of medieval Christendom in Europe, beliefs placed a high value on the knowledge of how men and women might achieve eternal salvation. {[}\ldots{]} That power and \textbf{authority} was reinforced by \textbf{control} over the means of communication, in the form of sacred books and of literacy in a common sacred language, Latin. & Means of communication \\
\midrule

p.125 & Aided by differences of language, national governments could use technology to keep \textbf{control} by censorship, by monopoly or by restrictive licensing over national systems of education, over national newspapers and broadcasting and even over the publication of books and periodicals. Thus, in this new knowledge structure, the \textbf{authority} of the Church was displaced by the extended \textbf{authority} of the scientific state. & National systems of education / national newspapers and broadcasting / publication of books and periodicals \\
\midrule

p.128--129 & In the Manhattan Project, the US government brought together an international team of top physicists from various countries. But it reserved to itself exclusive \textbf{control} over their discovery of how to apply the principles of nuclear fission to warfare. & Discovery of how to apply the principles of nuclear fission to warfare \\
\midrule

p.134 & The politically important point about these communication systems is, of course, that the bank's head office becomes the gatekeeper, \textbf{controlling} access to the system. & Access to the system \\
\midrule

p.134 & Walter Wriston, former head of Citibank (now Citicorp), has even suggested that `banking today is information' (Wriston, 1986). Similarly, it is the \textbf{control} over, and access to, these global systems that also allows the great grain and commodity trading companies to enjoy such an oligopolistic position as compared with either the producers (the farmers) or the end-users. & Global systems \\
\midrule

p.135 & But what is important is that now even the systems reserved to the Pentagon depend on, and could not operate without, the technical know-how and co-operation of the major transnational corporations. The possibility of total \textbf{control} and monopoly by the state (outside the Soviet Union and China) has seemingly gone for good. & Systems reserved to the Pentagon \\
\midrule

p.155 & Like sea transport, the way in which the global air transport system is structured rests on a political fiction: the notion that a state `\textbf{controls}' the airspace above its territory, in the same way that, from the eighteenth century until the second half of the twentieth century, it notionally \textbf{controlled} its `territorial waters'. {[}\ldots{]} The notion that any government can \textbf{control} what goes on in the air above it is obviously even more of a fiction than the notion that it can \textbf{control} what goes on in the seas around it. & Airspace / territorial waters \\
\midrule

p.188 & The United States has dominated and directed the negotiations in the GATT, {[}\ldots{]} partly because of the bargaining power conferred on it by its \textbf{control} over so large and rich a domestic market. & Domestic market \\
\midrule

p.193 & But the reason why the European Coal and Steel Community no longer held centre stage by the 1970s was much more economic than political. {[}\ldots{]} all the European states no longer \textbf{controlled} within their own frontiers their chief source of industrial energy, but were all in the same boat as net importers of oil and gas. & Source of industrial energy \\
\midrule

p.202 & The Iranian revolution of 1978--9, tempted the OPEC members to try the same gambit again {[}\ldots{]} OPEC found itself obliged to agree on a climb-down, a \$5 reduction in oil prices by which it hoped to show it was still in \textbf{control}. But, this time, the market took charge. & Oil market \\
\midrule

p.202 & The companies still had \textbf{control} of the technology of exploration, of offshore production, of refining and marketing; and they had the capital necessary for risk-taking in an essentially risky business. & Technology of exploration, offshore production, refining, marketing, capital \\
\midrule

p.237 & {[}\ldots{]} the United States has not in fact lost power in the world market economy. As that economy has grown and spread, the source of its power has shifted from the land and the people into \textbf{control} over structures of the world system. & Structures of the world system \\
\midrule

p.242 & {[}\ldots{]} the Americans reserved to themselves the right to decide unilaterally when to abandon sanctions and to use force against Saddam Hussein, and then to decide on the conduct of Desert Storm, and on how and when to call off the attack. Confidence in US leadership has also been undermined, not by a loss of power over others, but by lost \textbf{control} over its own tangled web of overblown bureaucracy. & Bureaucracy \\
\bottomrule
\end{longtable}

\noindent Source: STRANGE, 1988.

\normalsize

\section*{Appendix 4 -- Control, 1996}

Excerpts on \textbf{control} in \emph{The Retreat of the State: Diffusion of Power in the World Economy} by Susan Strange, 1996.

\small
\begin{longtable}{p{1.2cm} p{9.5cm} p{3.2cm}}
\toprule
\textbf{Page} & \textbf{Quote} & \textbf{Object} \\
\midrule
\endhead

p.100 & At the peak of their power over society, states claimed, and exercised, the right to \textbf{control} the substance of information --- by censorship, for example, of books or the press --- and to \textbf{control} the means by which information was communicated --- post, telegraph and telephone. & Substance of information / post, telegraph and telephone \\
\midrule

p.100 & Other governments have been forced by a combination of technological and economic change to give up their exclusive \textbf{control} for the sake of maintaining the competitiveness in world markets of the national economies for whose welfare they are held responsible. & Substance of information / post, telegraph and telephone \\
\midrule

p.105--106 & The power of governments which, for social policy reasons, might want to keep rural areas and lonely old people fully integrated into the communications system at minimal costs has clearly diminished. So has the \textbf{control} of governments. By means of their ownership of state monopolies, PTTs used to have \textbf{control} over the design and availability of such communications. No longer. {[}\ldots{]} Even the government of a country with a potential market as large as that of China no longer has the option of \textbf{controlling} and running its own communication system; the range of options open to it has narrowed to picking the foreign partners and negotiating with them the best terms of the alliance. & Post, telegraph and telephone (PTTs) / design, availability of PTTs communications \\
\midrule

p.108 & In telecommunications, the balance of benefit in the last two decades of the twentieth century would seem to have gone to the private sector firms at the expense of governments and their publicly owned and \textbf{controlled} enterprises. & Telecommunications enterprises \\
\midrule

p.112 & One estimate cited in 1995 by the chief prosecutor of Florence suggested that organised criminal groups in Russia then \textbf{controlled} 35 per cent of the commercial banks, 40 per cent of former State-owned industry, 35 per cent of private enterprise --- and as much as 60 per cent of commerce and 80 per cent of joint ventures with foreign firms. & Commercial banks, former state-owned industry, private enterprise, commerce, joint ventures with foreign firms \\
\midrule

p.120--121 & To reduce or even limit the economic wealth and potential for political and social disruption of these transnational criminal groups to manageable levels would strike at the very heart of national sovereignty {[}\ldots{]} If then suppression as an option is blocked by the refusal of state governments to give up their \textbf{control} of law enforcement, an alternative option would be to decriminalise the drug trade. & Law enforcement \\
\midrule

p.130 & By the early 1990s, it had run into balance of payments problems and needed the blessing and support of the IMF --- where, again, the US exercised the \textbf{controlling} veto power. & Veto power \\
\midrule

p.138 & But why, the political economist will ask, did states allow such great \textbf{authority} and influence to be exercised within the limits of the law by such a small number of private firms? One answer lies in the preference of governments in the Anglo-Saxon tradition for indirect rule, leaving it, wherever possible, to the operators themselves to monitor and \textbf{control} themselves, whether they are doctors, lawyers, stockbrokers or insurers. & Indirect rule within the law \\
\midrule

p.147 & There is also an element of private protectionism when a firm has monopoly \textbf{control} over a technology, or a system of marketing, or a brand-name that keeps away competitors. & Technology \\
\midrule

p.149 & Yet she admits that the most effective cartel of the four she studied, that in diamonds, almost entirely owed its success to the tight \textbf{control} over supplies exercised by one firm, Anglo-American, and the majority owners, the Oppenheimer family. & Supplies \\
\midrule

p.179 & Transfer of \textbf{authority} from the member governments to federal institutions over this central responsibility of political \textbf{authority} in a market economy has not yet happened. Nor has the replacement of national defence forces by a European Army under \textbf{control} of a European Chief of Staff. {[}\ldots{]} It seems that European governments --- though they are reluctant to say so --- really prefer a vacuum of power over key matters of security, currency, law and order and foreign policy to a real transfer of power to supranational institutions. & National defence forces \\
\midrule

p.187 & Other states which formerly had \textbf{controlled} and managed their national markets now found their PTTs challenged by the combination of new technology and foreign competitors. State policies changed in response. & National markets \\
\midrule

p.193 & There is no world central bank to exercise the judicious \textbf{control} over the creation of credit and the expansion of the money supply that historical experience has shown needs to be exercised. There is no world central bank with powers to \textbf{control} and regulate a banking system that operates transnationally in internationally integrated financial markets. & Banking system that operates transnationally in internationally integrated financial markets \\
\midrule

p.197 & The concentration of power in what Perry Anderson called the absolutist state was the means by which a politically \textbf{controlled} framework of rules was put round the emerging capitalist or market economy. & Political framework of rules \\
\bottomrule
\end{longtable}

\noindent Source: STRANGE, 1996.

\normalsize
